% Title pages: one page for each title layout of tikzflora, after the
% engraved and hand-coloured title pages of nineteenth-century flower
% books, with their display lines in \floraletters, which draws engraved
% capitals from the font's own outlines. Compile with LuaLaTeX:
%
%   lualatex titles      (or: latexmk titles)
%
% Each page sets its layout with \florathispage and marks its title with
% the floratitle environment; the guide on https://flora.alpacasoft.dev
% describes every key.
\documentclass[11pt]{article}
\usepackage[a4paper,margin=2.5cm]{geometry}
\usepackage{tikzflora}

\pagestyle{empty}
\setlength\parindent{0pt}
\florasetup{seed=5}

\begin{document}

% --- Flank ---------------------------------------------------------------
% Two plants rise up the sides of the page and meet at its foot, round a
% piece in the middle. The title sits in the open column between them.
\florathispage{layout=flank}
\vspace*{1cm}
\begin{floratitle}
\centering
{\fontsize{46}{52}\selectfont\floraletters[style=shaded]{FLORA'S}\par}
\vspace{.4em}
% arch bends the line through 24 degrees, as an engraver lettered it.
{\fontsize{30}{36}\selectfont\floraletters[style=engraved, arch=24]{DICTIONARY}\par}
\vspace{1em}
{\large\itshape of the flowers of the field\par}
\end{floratitle}
\vspace{1.5em}
{\centering\small\scshape Embellished in fountain-pen inks\par}
\vfill
{\centering\footnotesize\scshape Baltimore\par}
\clearpage

% --- Wreath --------------------------------------------------------------
% An open oval garland round the title.
\florathispage{layout=wreath}
\vspace*{\fill}
\begin{floratitle}
\centering
{\Large\floraletters[style=outline]{THE}\par}
\vspace{1.4em}
{\fontsize{34}{40}\selectfont\floraletters[style=shaded]{LANGUAGE}\par}
\vspace{1.2em}
{\Large\floraletters[style=outline]{OF}\par}
\vspace{1.2em}
{\fontsize{34}{40}\selectfont\floraletters[style=shaded]{FLOWERS.}\par}
\end{floratitle}
\vspace*{\fill}
{\centering\footnotesize\scshape London, at the sign of the acanthus\par}
\clearpage

% --- Bouquet -------------------------------------------------------------
% One upright bouquet; each line of the title sits on a ribbon wound
% through it.
\florathispage{layout=bouquet}
\vspace*{\fill}
\begin{floratitle}
\centering
\hspace*{3em}{\LARGE\floraletters[style=outline]{THE}\par}
\vspace{2.2em}
\hspace*{-2em}{\fontsize{30}{36}\selectfont\floraletters[style=engraved]{GARDEN}\par}
\vspace{2.2em}
\hspace*{1em}{\LARGE\floraletters[style=outline]{OF}\par}
\vspace{2.2em}
{\fontsize{30}{36}\selectfont\floraletters[style=engraved]{LANGUAGE}\par}
\end{floratitle}
\vfill
{\centering\footnotesize\scshape Fisher, Son \& Co.\ London\par}
\clearpage

% --- Scrollwork ----------------------------------------------------------
% An oval of gilt scrollwork round the title, with the flora grown over it.
\florathispage{layout=scrollwork}
\vspace*{\fill}
\begin{floratitle}
\centering
{\fontsize{28}{34}\selectfont\floraletters[style=engraved, arch=32, ink=tsuki-yo]{MY FIRST GARDEN}\par}
\vspace{.5em}
{\small\scshape A book for youth\par}
\vspace{2em}
{\footnotesize\itshape Consider the lilies of the field, how they grow;\\
they toil not, neither do they spin.\par}
\vspace{2em}
{\small by\par}
{\Large\floraletters[style=shaded]{A GARDENER}\par}
\end{floratitle}
\vspace*{\fill}
\end{document}
